% !Mode:: "TeX:UTF-8"
\chapter*{结论\markboth{结论}{}}
\addcontentsline{toc}{chapter}{结论}

本文针对模型未知的离散线性时不变系统，研究了以最少输入输出数据实现数据驱动范数最优迭代学习控制的方法，沿离线数据生成、控制算法设计、收敛性分析和仿真验证四个环节构建了一套完整方案。

在数据生成方面，本文提供了一种用于迭代学习控制的最短实验设计。具体而言，通过确保每个时刻新施加的输入避开一个低维仿射集，联合 Hankel 矩阵的秩便可严格增加 $1$ 直至达到理论上界 $Lm+n$。据此，以逐次随机输入的方式即可在理论最短数据长度 $(m+1)L+n-1$ 步内获得信息量充分的离线数据，远少于传统持续激励条件的要求。

在控制设计方面，本文基于 Willems' 基本引理，提出了一种数据驱动范数优化迭代学习控制策略。具体而言，当联合Hankel矩阵达到目标秩后，其列空间张成系统的全部行为，通过零初始条件响应矩阵 $W_0$ 将系统输入输出关系完全用数据矩阵替代，从而将范数最优迭代学习控制转化为仅依赖数据的凸优化问题，导出了控制更新律的解析解，并证明了该方法与基于模型的范数最优迭代学习控制等价。

此外，本文提供了针对数据驱动范数优化迭代学习控制的收敛性分析，严格证明了当联合 Hankel 矩阵达到目标秩时，跟踪误差范数沿迭代轴单调收敛至零，且当伴随系统正定时具有 $1/(1+\sigma^2)$ 的指数收敛速率，与基于模型方法的理论保证完全一致。

本文还提供了仿真结果用于验证联合 Hankel 矩阵秩随数据长度逐步增长的过程与理论预测吻合，最短数据长度与理论最小值一致；在两种不同期望轨迹下，所提出的数据驱动范数最优迭代学习控制方法均实现了高精度轨迹跟踪目标。

本文方法目前要求系统为线性时不变且每次迭代初始状态精确重置，尚未考虑量测噪声与实际系统验证。在后续研究中，将本文方法推广至时变系统是一个值得关注的方向\upcite{barton2011}。现实世界中大量被控对象本质上具有时变特征——工业机器人关节摩擦随温度和磨损渐变，飞行器气动参数随飞行高度和速度改变，化工过程反应速率受催化剂活性衰减影响。对于这类系统，如何设计具有在线数据更新能力的迭代学习控制策略，或建立时变系统的短数据实验准则，是兼具理论价值与工程意义的课题。此外，向非线性系统扩展、在噪声环境下建立鲁棒方法以及在实际机电平台上开展实验验证，也是后续工作的重要方向。
